\section*{第 \ref{ch:b2:blood-circulation} 章——血液与循环系统}

\begin{solution}{exo:b2:blood-circulation:1}
\omterm{def:b2:blood-circulation:blood}{血浆}占 \qty{55}{\%}（水、蛋白质 \qty{70}{g/L}、盐类、
营养物质、气体、激素）；细胞占 \qty{45}{\%}。\omterm{def:b2:blood-circulation:blood}{红细胞} $5\times
10^{12}$/L，直径 \qty{7}{\micro m} 的圆盘，负责氧运输；白细胞
$7\times 10^{9}$/L，直径 \qtyrange{10}{15}{\micro m}，负责免疫；\omterm{def:b2:blood-circulation:blood}{血小板}
$3\times 10^{11}$/L，\qty{2}{\micro m} 的碎片，参与\omterm{prop:b2:blood-circulation:clotting}{止血}。
\end{solution}

\begin{solution}{exo:b2:blood-circulation:2}
鱼类：心脏 $\to$ 鳃 $\to$ 身体 $\to$ 心脏，单循环，
身体所接收的是流经鳃后残留的低压\omterm{def:b2:blood-circulation:blood}{血液}。哺乳动物：
右心 $\to$ 肺（\qty{25}{mmHg}）$\to$ 左心 $\to$ 身体
（\qty{120}{mmHg}）$\to$ 右心。\omterm{def:b2:blood-circulation:circuits}{双循环}
使身体获得高压，而不使肺承受此压力，并使
两条循环可以分别调节。
\end{solution}

\begin{solution}{exo:b2:blood-circulation:3}
\omterm{def:b2:blood-circulation:vessels}{动脉}：管壁厚，有弹性层和肌肉——在高压下输送\omterm{def:b2:blood-circulation:blood}{血液}，
并使脉搏趋于平稳。\omterm{def:b2:blood-circulation:vessels}{小动脉}：管壁主要为
平滑肌——决定阻力并分配血流。\omterm{def:b2:blood-circulation:vessels}{毛细血管}：
管壁仅一个内皮细胞厚——负责交换。\omterm{def:b2:blood-circulation:vessels}{静脉}：管壁薄、可扩张，
有瓣膜——在低压下回流\omterm{def:b2:blood-circulation:blood}{血液}，并储存
其中大部分\omterm{def:b2:blood-circulation:blood}{血液}。
\end{solution}

\begin{solution}{exo:b2:blood-circulation:4}
$70\times 72 = \qty{5}{L/min}$，一天 \qty{7200}{L}，与
\qty{5}{L} 的血量相比：全部血量每天流经心脏 \num{1440}
次。没有任何器官能每天制造或消耗七吨\omterm{def:b2:blood-circulation:blood}{血液}，
因此同一批\omterm{def:b2:blood-circulation:blood}{血液}必须回流——即\omterm{def:b2:blood-circulation:blood}{血液}是循环的。
\end{solution}

\begin{solution}{exo:b2:blood-circulation:5}
$\Delta P = 8\eta LQ/\pi r^{4} = 8\times 3\times 10^{-3}\times
0.4\times 8.3\times 10^{-5}/(\pi\times 2.44\times 10^{-8}) =
\qty{10}{Pa}$，\qty{0.08}{mmHg}：主\omterm{def:b2:blood-circulation:vessels}{动脉}的损耗可忽略；
压降主要消耗在\omterm{def:b2:blood-circulation:vessels}{小动脉}中。
\end{solution}

\begin{solution}{exo:b2:blood-circulation:6}
$(15/12)^{4} = 2.4$：阻力升高为 2.4 倍，血流降至
\qty{41}{\%}。$(15/20)^{4} = 0.32$：阻力降至三分之一，血流
增至 3.2 倍。
\end{solution}

\begin{solution}{exo:b2:blood-circulation:7}
主\omterm{def:b2:blood-circulation:vessels}{动脉} $83/3 = \qty{28}{cm/s}$；\omterm{def:b2:blood-circulation:vessels}{毛细血管} $83/3000 =
\qty{0.028}{cm/s}$；通过时间 $0.1/0.028 = \qty{3.6}{s}$。运动时
心输出量升至五倍，开放\omterm{def:b2:blood-circulation:vessels}{毛细血管}的面积或许
增至三倍，因此通过时间降至约 \qty{2}{s}——仍
足以使氧气扩散出去。
\end{solution}

\begin{solution}{exo:b2:blood-circulation:8}
\omterm{def:b2:blood-circulation:vessels}{动脉}端 $35 - 25 = \qty{+10}{mmHg}$（滤过）；\omterm{def:b2:blood-circulation:vessels}{静脉}端 $15
- 25 = \qty{-10}{mmHg}$（重吸收）。若\omterm{def:b2:blood-circulation:vessels}{静脉}端 $P_c = 28$，
则净值为 $+3$：液体沿整个长度滤过，没有任何
液体被重吸收——即\omterm{prop:b2:blood-circulation:exchange}{水肿}，也就是心力衰竭时的脚踝肿胀。
\end{solution}

\begin{solution}{exo:b2:blood-circulation:9}
$40/66000 = \qty{0.61}{mol/m^3}$; $\pi = 0.61\times 8.314\times 310 =
\qty{1570}{Pa} = \qty{11.8}{mmHg}$；\omterm{thm:b2:blood-circulation:oncotic}{白蛋白}减半时为 \qty{5.9}{mmHg}。
钠可自由穿过\omterm{def:b2:blood-circulation:vessels}{毛细血管}壁，故其浓度在
两侧相等，因而不产生跨壁的渗透压。
\end{solution}

\begin{solution}{exo:b2:blood-circulation:10}
$RC = \qty{2}{s}$：$120e^{-0.4} = \qty{80}{mmHg}$。$C = 1$：$RC = 1$，
$120e^{-0.8} = \qty{54}{mmHg}$——同样的每搏输出量在僵硬的主\omterm{def:b2:blood-circulation:vessels}{动脉}中使
收缩压升高得更多。\omterm{thm:b2:blood-circulation:windkessel}{脉压}增大
意味着心脏射血时需对抗更高的峰值，做功更多，
同样输出量下耗氧也更多。
\end{solution}

\begin{solution}{exo:b2:blood-circulation:11}
静息：$90/83 = \qty{1.08}{mmHg.s/mL}$；运动：$105/417 =
\qty{0.25}{mmHg.s/mL}$，阻力降为原来的四分之一。由于 $R \propto r^{-4}$，
$r$ 增大为原来的 $4^{1/4} = 1.41$ 倍：\omterm{def:b2:blood-circulation:vessels}{小动脉}
平均扩张了 \qty{41}{\%}。
\end{solution}

\begin{solution}{exo:b2:blood-circulation:12}
连续性使各处流速由总截面决定，
而血管树的构造使\omterm{def:b2:blood-circulation:vessels}{毛细血管}处的截面
约为主\omterm{def:b2:blood-circulation:vessels}{动脉}的一千倍，在
\omterm{def:b2:blood-circulation:vessels}{静脉}中又变小；泊肃叶定律将阻力置于\omterm{def:b2:blood-circulation:vessels}{小动脉}，
那里的肌肉可以改变其阻力，并使粗大血管几乎
无损失。\omterm{def:b2:blood-circulation:circuits}{开放式循环}没有\omterm{def:b2:blood-circulation:vessels}{毛细血管}使\omterm{def:b2:blood-circulation:blood}{血液}在某一
确定位置减速，没有可以设定阻力的血管，也没有
办法将血流导向某一器官——它只能搅动。
\end{solution}

\begin{solution}{pb:b2:blood-circulation:1}
\textbf{1.} $70\times 72 = \qty{5.0}{L/min}$；一天 \qty{7300}{L}。
\textbf{2.} $7300/5 \approx 1450$ 次。
\textbf{3.} \qty{57}{g} 的四分之一，以每分钟 72 次搏动计：\qty{14}{g}
$\times 4320 = \qty{62}{kg}$ 每小时——约相当于一个成年男子的体重。
\textbf{4.} 没有任何食物能供应、也没有任何组织能消耗
相当于一个人体重的\omterm{def:b2:blood-circulation:blood}{血液}，每小时；唯一的出路是
同一批\omterm{def:b2:blood-circulation:blood}{血液}回流到心脏。
\textbf{5.} 系在手臂上的绳子使其下方的\omterm{def:b2:blood-circulation:vessels}{静脉}胀大，而非上方，
因此\omterm{def:b2:blood-circulation:vessels}{静脉}血流向心脏；将\omterm{def:b2:blood-circulation:vessels}{静脉}\omterm{def:b2:blood-circulation:blood}{血液}向手部挤压时，
\omterm{def:b2:blood-circulation:blood}{血液}会停在瓣膜处，无法被
推回——瓣膜只允许\omterm{def:b2:blood-circulation:blood}{血液}流向心脏。
\textbf{6.} 连接\omterm{def:b2:blood-circulation:vessels}{动脉}与\omterm{def:b2:blood-circulation:vessels}{静脉}的是\omterm{def:b2:blood-circulation:vessels}{毛细血管}；马尔皮基
于 1661 年在蛙肺中观察到了它们。
\textbf{7.} 面积 $\pi\times 1.25^{2} = \qty{4.9}{cm^2}$；$83/4.9 =
\qty{17}{cm/s}$.
\textbf{8.} $8\times 3\times 10^{-3}\times 0.4\times 8.3\times
10^{-5}/(\pi\times 2.44\times 10^{-8}) = \qty{10}{Pa}$, \qty{0.08}{mmHg}.
\textbf{9.} $R_{1} = 8\times 3\times 10^{-3}\times 10^{-3}/(\pi\times
5.06\times 10^{-20}) = \qty{1.5e14}{Pa.s/m^3}$；并联时，
$1.5\times 10^{14}/3\times 10^{6} = \qty{5.0e7}{Pa.s/m^3}$.
\textbf{10.} $8.3\times 10^{-5}\times 5.0\times 10^{7} = \qty{4200}{Pa}
= \qty{31}{mmHg}$.
\textbf{11.} 开放的\omterm{def:b2:blood-circulation:vessels}{毛细血管} $10^{10}$ 条，每条 $\pi(3\times 10^{-6})^{2}
= \qty{2.8e-11}{m^2}$: \qty{0.28}{m^2} $= \qty{2800}{cm^2}$; $v =
83/2800 = \qty{0.03}{cm/s}$；通过时间 $0.1/0.03 = \qty{3.3}{s}$。
\textbf{12.} 阻力 $\times(1/0.9)^{4} = 1.52$：\qty{47}{mmHg}；
心脏须将\omterm{def:b2:blood-circulation:vessels}{动脉}压提高 \qty{16}{mmHg}，否则
输出量将下降三分之一。
\textbf{13.} $+10$、$-10$ 和 $0$ mmHg。
\textbf{14.} 滤过量 $10\times 2 = \qty{20}{L}$ 每天；重吸收
$8.5\times 2 = \qty{17}{L}$.
\textbf{15.} 每天 \qty{3}{L} \omterm{prop:b2:blood-circulation:exchange}{淋巴}液。
\textbf{16.} $\pi_c = \qty{12.5}{mmHg}$：\omterm{def:b2:blood-circulation:vessels}{动脉}端净值为 $+22.5$
，\omterm{def:b2:blood-circulation:vessels}{静脉}端为 $+2.5$——全程滤过，没有
重吸收：\omterm{prop:b2:blood-circulation:exchange}{水肿}。
\textbf{17.} 净值 $+10$ 与 $+3$，平均 $+6.5$：滤过沿
整条\omterm{def:b2:blood-circulation:vessels}{毛细血管}进行，每天约 \qty{26}{L}；\omterm{prop:b2:blood-circulation:exchange}{淋巴}管无法
运走，液体便在组织中积聚，先从脚部开始。
\textbf{18.} $4\times 10^{10}\times 2\pi\times 3\times 10^{-6}\times
10^{-3} = \qty{750}{m^2}$; $t = x^{2}/2D = (2\times 10^{-5})^{2}/
(2\times 10^{-9}) = \qty{0.2}{s}$.
\textbf{19.} $(93 - 3)/83 = \qty{1.08}{mmHg.s/mL}$.
\textbf{20.} $RC = \qty{2.2}{s}$; $120e^{-0.8/2.2} = 120\times 0.69 =
\qty{83}{mmHg}$.
\textbf{21.} $RC = \qty{1.1}{s}$；$120e^{-0.73} = \qty{58}{mmHg}$：
\omterm{thm:b2:blood-circulation:windkessel}{脉压}由 37 增大到 \qty{62}{mmHg}。
\textbf{22.} $120e^{-0.3/2.2} = \qty{105}{mmHg}$：心率较快时
两次搏动之间压力几乎不下降。
\textbf{23.} $C\Delta P = 2\times 20 = \qty{40}{mL}$ 被储存；其余
\qty{30}{mL} 在收缩期内即经\omterm{def:b2:blood-circulation:vessels}{小动脉}流走。
\textbf{24.} 收缩的心室在收缩期挤闭自身的血管，
因此心肌在舒张期灌注，靠的是回缩的主\omterm{def:b2:blood-circulation:vessels}{动脉}所维持的
压力；若僵硬的主\omterm{def:b2:blood-circulation:vessels}{动脉}使舒张压塌陷，
心脏便会在两次搏动之间缺血。
\textbf{25.} 输出量 \qty{5}{L/min}；\omterm{def:b2:blood-circulation:vessels}{小动脉}压降 \qty{31}{mmHg}；
净滤液每天 \qty{3}{L} 进入\omterm{prop:b2:blood-circulation:exchange}{淋巴}；舒张压
当 $C = 2$ 时为 \qty{83}{mmHg}，当 $C = 1$ 时为 \qty{58}{mmHg}。
\end{solution}
